\scanpage{142}
1) \uline{Connexions sur torseurs affines}

$S/T$ schéma relatif \inserted{$(\unsure{g|q}:S\mathbin{-}T)$}, $E$ module qu. \unsure{ch.|coh.} sur $S$, d'où
\[
0\longrightarrow E\otimes\Omega^1_{S/T}\longrightarrow P^1_{S/T}(E)
\longrightarrow E\longrightarrow0 \tag*{(1)}
\]
extension de $\underline O_S$-modules. On a un splitting $g^{-1}(\underline O_T)$-linéaire
(mais pas $\underline O_S$-linéaire)
\[
d_E^1:E\longrightarrow P^1_{S/T}(E) \tag*{(2)}
\]
qui permet d'identifier $P^1_{S/T}(E)$ \uline{additiv$^{\mathrm{mt}}$} à $E\oplus E\otimes\Omega^1_{S/T}$, en
écrivant chaque \deleted{élément}\inserted{section} de $P^1_{S/T}(E)$
\inserted{(sur un ouvert)} de façon unique comme
\[
\zeta=\deleted{\unsure{\omega+}}\underset{\deleted{\unsure{E\otimes\Omega^1_{S/T}}}}{\,}
 d_E^{(1)}x+\omega\quad
 \left\{\begin{aligned}\omega&\in\Gamma(\ ,E\otimes\Omega^1_{S/T})\\
 x&\in\Gamma(\ ,E)\end{aligned}\right. \tag*{(3)}
\]
On a, si \deleted{\ensuremath{f}} $f$ est une section de $\underline O_S$
\[
f\zeta=f\omega+fd_E^{(1)}x=f\omega+
\underset{=\ d^1_{\unsure{\bullet}}(f)d^1_E(x)}{\underbrace{d_E^{(1)}(fx)}}+
\underset{-df}{[f-d_{\unsure{E|\bullet}}^{(1)}(f)]}d_E^1(x)
\]
i.e.
\[
\deleted{f(\omega+d_E^{(1)}x)=f\omega}\tag*{\deleted{(4)}}
\]
\[
f(d_E^1x+\omega)=d_E^1fx+(f\omega-x\otimes df).\tag*{(4)}
\]

Une \uline{connexion} \inserted{$\delta_E$} de $E$ rel. à $T$ est un splitting \inserted{$\underline O_S$-linéaire} de (1) ;
\deleted{\unsure{don|donc}} tant que splitting additif, compte tenu de
(2), il est donc donné par un hom additif
\[
d_E^0:E\longrightarrow E\otimes\Omega^1_{S/T} \tag*{(5)}
\]
par
\[
\delta_E x=d_E^1(x)-d_E^0(x), \tag*{(6)}
\]
il faut expliciter la condition sur $d_E^0$ de (5),
pour que $\delta_X$ défini par (6) soit $\underline O_S$-linéaire,
% archive sheet number: 141
